\documentclass[12pt,a4paper]{article}

% ── Packages ──────────────────────────────────────────────────────────────────
\usepackage[a4paper, margin=2.5cm]{geometry}
\usepackage{amsmath}
\usepackage{booktabs}
\usepackage{array}
\usepackage{tabularx}
\usepackage{xcolor}
\usepackage{listings}
\usepackage{algorithm}
\usepackage{algpseudocode}
\usepackage{hyperref}
\usepackage{graphicx}
\usepackage{caption}
\usepackage{subcaption}
\usepackage{tcolorbox}
\usepackage{titlesec}
\usepackage{parskip}
\usepackage{enumitem}
\usepackage{microtype}
\usepackage{fancyhdr}
\usepackage{mdframed}

% ── Colours ───────────────────────────────────────────────────────────────────
\definecolor{darkblue}{RGB}{0,70,127}
\definecolor{lightblue}{RGB}{220,235,247}
\definecolor{codebg}{RGB}{245,245,245}
\definecolor{codeframe}{RGB}{180,180,180}
\definecolor{boxbg}{RGB}{240,248,255}

% ── Hyperref ──────────────────────────────────────────────────────────────────
\hypersetup{
    colorlinks=true,
    linkcolor=darkblue,
    urlcolor=darkblue,
    citecolor=darkblue
}

% ── Section title formatting ──────────────────────────────────────────────────
\titleformat{\section}{\large\bfseries\color{darkblue}}{}{0em}{}[\titlerule]
\titleformat{\subsection}{\normalsize\bfseries\color{darkblue}}{\thesubsection}{1em}{}
\titleformat{\subsubsection}{\normalsize\itshape}{\thesubsubsection}{1em}{}

% ── Header / Footer ───────────────────────────────────────────────────────────
\pagestyle{fancy}
\fancyhf{}
\lhead{\small DAS 839 -- NoSQL Systems}
\rhead{\small Assignment II Report}
\cfoot{\thepage}
\renewcommand{\headrulewidth}{0.4pt}

% ── Code listing style ────────────────────────────────────────────────────────
\lstdefinestyle{pseudostyle}{
    backgroundcolor=\color{codebg},
    frame=single,
    rulecolor=\color{codeframe},
    basicstyle=\ttfamily\small,
    keywordstyle=\bfseries,
    breaklines=true,
    tabsize=4,
    xleftmargin=0.5cm,
    xrightmargin=0.5cm,
    aboveskip=6pt,
    belowskip=6pt,
}
\lstdefinestyle{shellstyle}{
    backgroundcolor=\color{black},
    frame=single,
    rulecolor=\color{codeframe},
    basicstyle=\ttfamily\small\color{white},
    breaklines=true,
    xleftmargin=0.5cm,
    xrightmargin=0.5cm,
    aboveskip=6pt,
    belowskip=6pt,
}

% ── tcolorbox styles ──────────────────────────────────────────────────────────
\tcbuselibrary{skins,breakable}
\newtcolorbox{objectivebox}{
    colback=lightblue,
    colframe=darkblue,
    fonttitle=\bfseries,
    title=Objective,
    breakable
}
\newtcolorbox{insightbox}[1][]{
    colback=boxbg,
    colframe=darkblue!60,
    fontupper=\small,
    left=6pt, right=6pt, top=6pt, bottom=6pt,
    #1,
    breakable
}

% ── Algorithm style ───────────────────────────────────────────────────────────
\algrenewcommand\algorithmicrequire{\textbf{Input:}}
\algrenewcommand\algorithmicensure{\textbf{Output:}}

% ─────────────────────────────────────────────────────────────────────────────
\begin{document}

% ── Cover Page ────────────────────────────────────────────────────────────────
\begin{titlepage}
    \centering
    \vspace*{1cm}
    {\color{darkblue}\rule{\linewidth}{1.5pt}}\\[0.4cm]
    {\LARGE\bfseries\color{darkblue} DAS 839 -- NoSQL Systems}\\[0.3cm]
    {\large Assignment II: MapReduce \& Apache Hadoop}\\[0.2cm]
    {\color{darkblue}\rule{\linewidth}{1.5pt}}\\[2cm]

    \begin{tabular}{rl}
        \textbf{Course}         & DAS 839 -- NoSQL Systems \\[0.4cm]
        \textbf{Group Members}  & \textit{[Your Name -- Roll Number]} \\
                                & \textit{[Group Member 2 -- Roll Number]} \\
                                & \textit{[Group Member 3 -- Roll Number]} \\[0.4cm]
        \textbf{GitHub Repo}    & \textit{[Insert GitHub URL here]} \\[0.4cm]
        \textbf{Submitted}      & April 14, 2026 \\
    \end{tabular}

    \vfill

    \begin{tcolorbox}[colback=lightblue, colframe=darkblue, width=0.95\textwidth]
        \textbf{Submission Notes:} Run instructions are provided in
        \texttt{problem\_2.txt}. All experiments were conducted on the Wikipedia
        English dump \texttt{Wikipedia-EN-20120601\_ARTICLES} (10,000 articles).
        The full pipeline is automated via \texttt{scripts/run\_all.sh}.
    \end{tcolorbox}
\end{titlepage}

% ── Table of Contents ─────────────────────────────────────────────────────────
\tableofcontents
\newpage

% =============================================================================
\section{Problem 2: Indexing Documents via Hadoop}
% =============================================================================

\begin{objectivebox}
Compute Document Frequency (DF) for all terms in the Wikipedia corpus
(\texttt{Wikipedia-EN-20120601\_ARTICLES}), identify the top-100 high-DF terms,
then compute TF-IDF scores for those terms across all documents using a
\emph{Stripes} MapReduce approach.
The Porter Stemmer from \texttt{opennlp-tools-1.9.3.jar} is used for term
normalisation. \texttt{CombineTextInputFormat} is used in both jobs to merge
10,000 small input files into large splits, drastically reducing map-task
launch overhead.
\end{objectivebox}

% -----------------------------------------------------------------------------
\subsection{Part (a): Document Frequency Computation}
% -----------------------------------------------------------------------------

\subsubsection{Pseudocode}

\begin{algorithm}[H]
\caption{Document Frequency -- Map and Reduce}
\begin{algorithmic}[1]
\State \textbf{// Setup: load stopwords and initialise stemmer}
\Function{setup}{context}
    \State $\text{stopwords} \leftarrow \text{load\_from\_distributed\_cache}()$
    \State $\text{stemmer} \leftarrow \textbf{new } \text{PorterStemmer}()$
\EndFunction

\State
\State \textbf{// Mapper}
\Function{Map}{key, line}
    \State $\text{docId} \leftarrow \text{extract\_filename}(\texttt{mapreduce.map.input.file})$
    \State $\text{tokens} \leftarrow \text{tokenize}(\text{line.toLowerCase}(), \; \text{split on } [\hat{\;}a\text{-}z]+)$
    \State $\text{seenThisLine} \leftarrow \{\}$
    \For{each $\text{token}$ in $\text{tokens}$}
        \If{$\neg \text{isValidToken}(\text{token})$} \textbf{continue} \EndIf
        \If{$\text{token} \in \text{stopwords}$} \textbf{continue} \EndIf
        \State $\text{stemmed} \leftarrow \text{stemmer.stem}(\text{token})$
        \If{$\neg \text{isValidToken}(\text{stemmed})$} \textbf{continue} \EndIf
        \If{$\text{stemmed} \in \text{stopwords}$} \textbf{continue} \EndIf
        \If{$\text{stemmed} \notin \text{seenThisLine}$}
            \State $\text{seenThisLine.add}(\text{stemmed})$
            \State $\text{emit}(\text{stemmed},\; \text{docId})$
        \EndIf
    \EndFor
\EndFunction

\State
\State \textbf{// Reducer}
\Function{Reduce}{term, docIds}
    \State $\text{distinctDocs} \leftarrow \{\}$
    \For{each $\text{docId}$ in $\text{docIds}$}
        \State $\text{distinctDocs.add}(\text{docId})$
    \EndFor
    \State $\text{emit}(\text{term},\; |\text{distinctDocs}|)$
\EndFunction
\end{algorithmic}
\end{algorithm}

\smallskip
\noindent\textit{where} \texttt{isValidToken(t)} returns \texttt{true} iff $t$
is purely alphabetic, has length $> 2$, and is not a Roman numeral
(\texttt{i}, \texttt{ii}, \ldots).

% ·············································································
\subsubsection{Program Logic}

The Document Frequency program counts the number of distinct Wikipedia
documents in which each unique stemmed term appears.

\begin{itemize}[leftmargin=1.4em]
    \item \textbf{Setup Phase:} The stopwords file is loaded from the
    \emph{Distributed Cache} into a \texttt{HashSet} inside the Mapper's
    \texttt{setup()} method.  A \texttt{PorterStemmer} instance is created
    once per mapper for term normalisation.

    \item \textbf{Input Splitting:} \texttt{CombineTextInputFormat} is
    configured with a 64\,MB maximum split size.  This consolidates 10,000
    individual article files into just \textbf{3 actual map tasks}, eliminating
    per-task JVM startup overhead.

    \item \textbf{Document ID Derivation:} Because
    \texttt{CombineTextInputFormat} packs multiple files into one split, the
    document ID cannot be read directly from \texttt{FileSplit.getPath()}.
    Instead, the mapper reads the \texttt{mapreduce.map.input.file}
    configuration property (set per-line by the framework), strips the file
    extension, and uses the resulting filename as the document ID.

    \item \textbf{Mapper:} Each input line is lowercased and split on
    non-alphabetic characters.  Tokens are filtered by \texttt{isValidToken()},
    checked against the stopword set, and stemmed using \texttt{PorterStemmer}.
    A per-line \texttt{HashSet} (\texttt{seenThisLine}) ensures each stemmed
    term is emitted at most once per line, enforcing document-level uniqueness.
    The mapper emits \texttt{(stemmed\_term, docId)}.

    \item \textbf{Reducer:} All incoming document IDs for a given term are
    collected into a \texttt{HashSet} to deduplicate across lines from different
    mappers.  The reducer emits \texttt{(term, distinctCount)}, the final DF
    value.
\end{itemize}

% ·············································································
\subsubsection{Post-processing for Top-100 Terms}

After the MapReduce job completes, the top-100 highest document-frequency
terms are extracted and uploaded to HDFS for use in Part (b):

\begin{lstlisting}[style=shellstyle]
hdfs dfs -cat /user/krish/output/df_output/part-r-* \
  | sort -t$'\t' -k2 -nr \
  | head -100 > df_top100.tsv

hdfs dfs -put -f df_top100.tsv /user/krish/friend_df_top100.tsv
\end{lstlisting}

% ·············································································
\subsubsection{Runtime Analysis}

\begin{lstlisting}[style=shellstyle]
Job completed in     : 57s
Map tasks launched   : 5  (3 data-local, 2 killed speculative)
Actual map splits    : 3
Reduce tasks         : 1
Map input records    : 10,000
Map output records   : 5,813,551
Reduce input groups  : 409,843
Reduce output records: 409,843
HDFS bytes read      : 167,289,868
HDFS bytes written   :   4,386,938
CPU time spent       : 160,990 ms
Peak Map Phys. Mem.  : 449 MB
Peak Reduce Phys. Mem: 412 MB
\end{lstlisting}

\begin{insightbox}
\textbf{CombineTextInputFormat benefit:}
Compressing 10,000 files into 3 splits eliminated the per-task JVM launch
overhead that would otherwise dominate runtime.

\textbf{Parallelism:}
The aggregate CPU across map and reduce slots ($\approx 154$\,s) exceeds the
57\,s wall-clock time, confirming effective parallel execution across cores.

\textbf{Shuffle overhead:}
Map output materialized bytes (90\,MB) exceed raw output bytes (78\,MB) due to
serialisation and two spill passes
(\texttt{Spilled Records = 11,627,102}), caused by the large volume of
intermediate \texttt{(term, docId)} pairs.

\textbf{System time is low:}
Minimal system-level overhead indicates that the bottleneck is CPU computation
(tokenisation, stemming, hashing) rather than I/O.
\end{insightbox}

% ·············································································
\subsubsection{Output Schema}

The output TSV file has the following schema:

\begin{center}
\texttt{TERM\textbackslash tDF}
\end{center}

where \texttt{TERM} is the Porter-stemmed token and \texttt{DF} is the count of
distinct Wikipedia articles containing it.

% ·············································································
\subsubsection{Top-100 High Document Frequency Terms}

Total unique stemmed terms in the corpus: \textbf{409,843}.
The following table shows the top-20 terms out of the 100 extracted.

\begin{table}[h]
\centering
\caption{Top-20 Terms by Document Frequency}
\begin{tabular}{@{}clr@{}}
\toprule
\textbf{Rank} & \textbf{Term (stemmed)} & \textbf{DF} \\
\midrule
1  & apo      & 9518 \\
2  & categori & 9381 \\
3  & refer    & 9340 \\
4  & quot     & 8674 \\
5  & other    & 8415 \\
6  & link     & 8296 \\
7  & thi      & 8190 \\
8  & extern   & 8008 \\
9  & two      & 7949 \\
10 & includ   & 7885 \\
11 & year     & 7801 \\
12 & first    & 7689 \\
13 & http     & 7418 \\
14 & time     & 7349 \\
15 & state    & 7299 \\
16 & name     & 6838 \\
17 & www      & 6804 \\
18 & see      & 6761 \\
19 & histori  & 6662 \\
20 & peopl    & 6626 \\
\bottomrule
\end{tabular}
\end{table}

% ── SCREENSHOT PLACEHOLDER ────────────────────────────────────────────────────
\begin{mdframed}[backgroundcolor=yellow!15, linecolor=orange, linewidth=1pt]
\textbf{[SCREENSHOT 1 -- Insert here]}
Terminal output showing the full top-100 list from \texttt{df\_top100.tsv}
(the block starting with \texttt{apo~9518} down to \texttt{home~4014} from
your run).  Use: \verb|cat df_top100.tsv| or the terminal block from your
\texttt{run\_all.sh} output.
\end{mdframed}

% ── SCREENSHOT PLACEHOLDER ────────────────────────────────────────────────────
\begin{mdframed}[backgroundcolor=yellow!15, linecolor=orange, linewidth=1pt]
\textbf{[SCREENSHOT 2 -- Insert here]}
YARN ResourceManager Web UI screenshot showing Job 2a
(\texttt{job\_1776036771067\_0005}) in SUCCEEDED state.
URL: \texttt{http://krish-Inspiron-15-3520:8088}
\end{mdframed}

% ── SCREENSHOT PLACEHOLDER ────────────────────────────────────────────────────
\begin{mdframed}[backgroundcolor=yellow!15, linecolor=orange, linewidth=1pt]
\textbf{[SCREENSHOT 3 -- Insert here]}
Full YARN job counters terminal output for Job 2a (the block beginning with
\texttt{File System Counters} through \texttt{Bytes Written=4386938}).
\end{mdframed}

% -----------------------------------------------------------------------------
\subsection{Part (b): TF-IDF Score Computation -- Stripes Approach}
% -----------------------------------------------------------------------------

\subsubsection{Pseudocode}

\begin{algorithm}[H]
\caption{TF-IDF Score -- Stripes MapReduce}
\begin{algorithmic}[1]
\State \textbf{// Setup (Mapper \& Reducer): load df\_top100.tsv from Distributed Cache}
\Function{setup}{context}
    \State $\text{dfMap} \leftarrow \text{load\_tsv\_as\_map}(\text{cache\_file})$
    \hfill \(\triangleright\) Maps term $\to$ DF
\EndFunction

\State
\State \textbf{// Mapper -- Stripes with in-mapper combining}
\Function{Map}{key, line}
    \State $\text{docId} \leftarrow \text{extract\_filename}(\texttt{mapreduce.map.input.file})$
    \State $\text{stripe} \leftarrow \text{docStripes}[\text{docId}]$
    \hfill \(\triangleright\) get or create in-memory \texttt{HashMap}
    \State $\text{tokens} \leftarrow \text{tokenize}(\text{line.toLowerCase}())$
    \For{each $\text{token}$ in $\text{tokens}$}
        \If{$\neg \text{isValidToken}(\text{token})$} \textbf{continue} \EndIf
        \State $\text{stemmed} \leftarrow \text{stemmer.stem}(\text{token})$
        \If{$\text{isValidToken}(\text{stemmed})$ \textbf{and} $\text{stemmed} \in \text{dfMap}$}
            \State $\text{stripe}[\text{stemmed}] \mathrel{+}= 1$
        \EndIf
    \EndFor
\EndFunction

\State
\Function{cleanup}{context}
    \For{each $(\text{docId}, \text{stripe})$ in $\text{docStripes}$}
        \State $\text{stripeStr} \leftarrow \text{serialize}(\text{stripe})$
        \hfill \(\triangleright\) \texttt{"t1:c1|t2:c2|..."}
        \State $\text{emit}(\text{docId},\; \text{stripeStr})$
    \EndFor
\EndFunction

\State
\State \textbf{// Reducer}
\Function{Reduce}{docId, stripeStrings}
    \State $\text{termCounts} \leftarrow \{\}$
    \For{each $\text{stripeStr}$ in $\text{stripeStrings}$}
        \State $\text{pairs} \leftarrow \text{split}(\text{stripeStr},\; \text{``|''})$
        \For{each pair}
            \State $(\text{term}, \text{count}) \leftarrow \text{split}(\text{pair}, \; \text{``:''})$
            \State $\text{termCounts}[\text{term}] \mathrel{+}= \text{parseInt}(\text{count})$
        \EndFor
    \EndFor
    \For{each $(\text{term}, \text{tf})$ in $\text{termCounts}$}
        \State $\text{df} \leftarrow \text{dfMap}[\text{term}]$
        \State $\text{score} \leftarrow \text{tf} \times \log(10000.0 / \text{df} + 1)$
        \State $\text{emit}(\text{docId},\; \text{term} + \texttt{"\textbackslash t"} + \text{format}(\text{score}, 6))$
    \EndFor
\EndFunction
\end{algorithmic}
\end{algorithm}

% ·············································································
\subsubsection{Program Logic}

The TF-IDF Scorer computes a relevance score for each (document, term) pair
restricted to the top-100 high-DF terms identified in Part~(a).

\begin{itemize}[leftmargin=1.4em]
    \item \textbf{Setup Phase (Mapper and Reducer):} Both the Mapper and
    Reducer load \texttt{df\_top100.tsv} from the \emph{Distributed Cache}
    into an in-memory \texttt{HashMap\textless String, Integer\textgreater}
    (\texttt{dfMap}).  Loading in the reducer ensures TF-IDF score computation
    does not require a second MapReduce pass.

    \item \textbf{Input Splitting:} Same \texttt{CombineTextInputFormat}
    strategy as Part~(a) -- 64\,MB max split, reducing 10,000 article files
    to \textbf{3 actual map tasks}.

    \item \textbf{Mapper -- Stripes with In-Mapper Combining:}
    The mapper does \emph{not} emit inside \texttt{map()}.  Instead, for each
    input line it reads the document ID from \texttt{mapreduce.map.input.file},
    stems each valid token, and if the stemmed token exists in \texttt{dfMap},
    increments its count in an in-memory \texttt{docStripes} map
    (\texttt{HashMap\textless docId, HashMap\textless term,
    count\textgreater\textgreater}).  This is the classic \emph{Stripes} design
    pattern: one stripe (associative array) per document is maintained
    throughout the mapper's lifetime.  Emission happens only in
    \texttt{cleanup()}, where each document's stripe is serialised and emitted
    as \texttt{(docId, stripeString)}.

    \item \textbf{Stripe Serialisation Format:} Each stripe is encoded as a
    pipe-delimited string of colon-separated \texttt{term:count} pairs, e.g.:
    \begin{center}
    \texttt{american:3|extern:2|http:1|apo:7|\ldots}
    \end{center}

    \item \textbf{Reducer:} For each document ID the reducer receives one or
    more partial stripe strings (one per map task that processed lines of that
    document).  It merges them by summing counts per term, computes the TF-IDF
    score, and emits one output row per (docId, term) pair.
\end{itemize}

% ·············································································
\subsubsection{Scoring Formula}

\begin{equation}
    \text{SCORE}(t,\, d)
    \;=\;
    \text{TF}(t, d)
    \;\times\;
    \log\!\left(\frac{10000}{\text{DF}(t)} + 1\right)
\end{equation}

\noindent where:
\begin{itemize}[leftmargin=1.4em]
    \item $\text{TF}(t, d)$ = number of occurrences of term $t$ in document $d$
    \item $\text{DF}(t)$ = number of distinct documents containing term $t$
    (computed in Part~a)
    \item $10000$ = total number of Wikipedia articles in the corpus
\end{itemize}

% ·············································································
\subsubsection{Runtime Analysis}

\begin{lstlisting}[style=shellstyle]
Job completed in      : 53s
Map tasks launched    : 5  (3 data-local, 2 killed speculative)
Actual map splits     : 3
Reduce tasks          : 1
Map input records     : 10,000
Map output records    : 10,000       <- one stripe per document (cleanup)
Map output bytes      : 4,391,675
Reduce input groups   : 10,000
Reduce output records : 541,705
HDFS bytes read       : 167,289,868
HDFS bytes written    :  11,707,160
CPU time spent        : 158,460 ms
Peak Map Phys. Mem.   : 521 MB
Peak Reduce Phys. Mem.: 339 MB
\end{lstlisting}

\begin{insightbox}
\textbf{In-mapper combining drastically reduces shuffle volume:}
Map output records = 10,000 (one per document) vs.\ 5,813,551 in Part~(a).
The stripes pattern aggregates all term counts in-memory before emitting, so
the shuffle phase transfers only $\approx 4.4$\,MB instead of $\approx 90$\,MB.

\textbf{Faster wall-clock time than Part~(a):}
53\,s vs.\ 57\,s, despite larger output (541,705 rows).  The savings come from
near-zero spill overhead (\texttt{Spilled Records = 20,000} vs.\ 11.6\,M)
and a much smaller shuffle.

\textbf{Higher peak mapper memory:}
The \texttt{docStripes} map accumulates TF counts for all documents assigned
to a mapper simultaneously ($\approx 3{,}333$ docs per split).  This raised
peak map memory to 521\,MB vs.\ 449\,MB in Part~(a) -- the expected trade-off
of in-mapper combining.

\textbf{Stripes efficiency:}
Restricting processing to only 100 terms (checked via
\texttt{dfMap.containsKey()}) means the mapper skips the vast majority of
tokens, keeping stripes small and the whole job lightweight despite a full
corpus scan.
\end{insightbox}

% ·············································································
\subsubsection{Output Schema}

The output TSV file has the following schema:

\begin{center}
\texttt{ID\textbackslash tTERM\textbackslash tSCORE}
\end{center}

where \texttt{ID} is the document filename without extension (e.g.,
\texttt{100005}), \texttt{TERM} is the Porter-stemmed token, and \texttt{SCORE}
is formatted to 6 decimal places.

% ·············································································
\subsubsection{Representative TF-IDF Output}

\begin{table}[h]
\centering
\caption{Sample TF-IDF Output Rows (Document \texttt{100005})}
\begin{tabular}{@{}llr@{}}
\toprule
\textbf{ID} & \textbf{TERM} & \textbf{SCORE} \\
\midrule
100005 & american & 1.035975 \\
100005 & extern   & 0.810375 \\
100005 & three    & 1.024390 \\
100005 & major    & 1.121527 \\
100005 & www      & 4.520531 \\
100005 & decemb   & 1.204089 \\
100005 & press    & 3.562563 \\
100005 & group    & 4.779429 \\
100005 & area     & 1.096882 \\
100005 & apo      & 7.899676 \\
100005 & august   & 2.429295 \\
100005 & work     & 2.089465 \\
100005 & gener    & 0.999107 \\
100005 & remain   & 1.242151 \\
100005 & develop  & 1.197756 \\
100005 & both     & 1.041519 \\
100005 & take     & 1.211492 \\
100005 & includ   & 0.819000 \\
100005 & name     & 0.901143 \\
100005 & http     & 11.096731 \\
\bottomrule
\end{tabular}
\end{table}

% ── SCREENSHOT PLACEHOLDER ────────────────────────────────────────────────────
\begin{mdframed}[backgroundcolor=yellow!15, linecolor=orange, linewidth=1pt]
\textbf{[SCREENSHOT 4 -- Insert here]}
Terminal output showing the sample TF-IDF rows from \texttt{tfidf\_output.tsv}
(the block starting with \texttt{100005~american~1.035975} from your
\texttt{run\_all.sh} output).
\end{mdframed}

% ── SCREENSHOT PLACEHOLDER ────────────────────────────────────────────────────
\begin{mdframed}[backgroundcolor=yellow!15, linecolor=orange, linewidth=1pt]
\textbf{[SCREENSHOT 5 -- Insert here]}
YARN ResourceManager Web UI screenshot showing Job 2b
(\texttt{job\_1776036771067\_0006}) in SUCCEEDED state.
URL: \texttt{http://krish-Inspiron-15-3520:8088}
\end{mdframed}

% ── SCREENSHOT PLACEHOLDER ────────────────────────────────────────────────────
\begin{mdframed}[backgroundcolor=yellow!15, linecolor=orange, linewidth=1pt]
\textbf{[SCREENSHOT 6 -- Insert here]}
Full YARN job counters terminal output for Job 2b (the block beginning with
\texttt{File System Counters} through \texttt{Bytes Written=11707160}).
\end{mdframed}

% =============================================================================
\newpage
\appendix
\section*{Appendix}
\addcontentsline{toc}{section}{Appendix}
% =============================================================================

\subsection*{A. System Configuration}

\begin{table}[h]
\centering
\caption{Experimental System Configuration}
\begin{tabular}{@{}ll@{}}
\toprule
\textbf{Parameter} & \textbf{Value} \\
\midrule
Hadoop Version   & 3.x (YARN pseudo-distributed mode) \\
Java Version     & Java 8 \\
Dataset          & Wikipedia-EN-20120601\_ARTICLES (10,000 files, $\approx$160\,MB) \\
Input Format     & \texttt{CombineTextInputFormat} (64\,MB max split) \\
Stemmer          & OpenNLP PorterStemmer (\texttt{opennlp-tools-1.9.3.jar}) \\
Stopwords        & \texttt{stopwords.txt} (loaded via Distributed Cache) \\
\bottomrule
\end{tabular}
\end{table}

\subsection*{B. Run Instructions}

Run instructions are provided in \texttt{problem\_2.txt}.
The pipeline is automated via a single script:

\begin{lstlisting}[style=shellstyle]
cd /path/to/code
./scripts/run_all.sh
\end{lstlisting}

This script:
\begin{enumerate}
    \item Builds the JAR with Maven (\texttt{mvn package -q})
    \item Runs \texttt{parta.DocumentFrequency} on the Wikipedia input
    \item Extracts and uploads \texttt{df\_top100.tsv} to HDFS
    \item Runs \texttt{partb.TFIDFScorer} using the top-100 DF file
\end{enumerate}

\subsection*{C. Tokenisation and Normalisation}

The tokenisation logic is consistent across both MapReduce jobs:

\begin{itemize}
    \item Each line is converted to lowercase and split using
    \texttt{[{\textasciicircum}a-z]+}, removing all non-alphabetic separators.
    \item Tokens are filtered to remove:
    \begin{itemize}
        \item tokens of length two or less,
        \item non-alphabetic tokens,
        \item Roman numerals (\texttt{i}, \texttt{ii}, \texttt{iii}, \ldots), and
        \item stopwords loaded from \texttt{stopwords.txt} via Distributed Cache.
    \end{itemize}
    \item Remaining tokens are stemmed using the OpenNLP Porter Stemmer.
    Post-stemming validity and stopword checks are also applied.
\end{itemize}

\end{document}
